%%%%%%%%%%%%%%%%%%%%%%%%%%%%%%%%%%%%%%%%%%%%%%%%%%%%%%%%%%%%%
%% Begin exercise %%
%%%%%%%%%%%%%%%%%%%%%%%%%%%%%%%%%%%%%%%%%%%%%%%%%%%%%%%%%%%%%
\ex{PLC}
% \newif%\ifshowsolutions
% \showsolutionsfalse % <-- change to \showsolutionstrue to show
%%%%%%%%%%%%%%%%%%%%%%%%%%%%%%%%%%%%%%%%%%%%%%%%%%%%%%%%%%%%%
%% Task 5.1 %%
%%%%%%%%%%%%%%%%%%%%%%%%%%%%%%%%%%%%%%%%%%%%%%%%%%%%%%%%%%%%%

\task{Logic Analysis (1)}

The logic below is given:
\begin{figure}[h!]
    \centering
    \includegraphics[width=0.65\textwidth]{task_5_1.pdf}
    % \caption{Logic Circuit}
\end{figure}
Please analyse it and simplfy it !

\subtask{Analog Output Voltage}

%\ifshowsolutions
\begin{solutionblock}

\begin{center}
\begin{tabular}{|>{\centering\arraybackslash}p{1cm}|
                >{\centering\arraybackslash}p{1cm}|
                >{\centering\arraybackslash}p{1cm}|}
\hline
A & B & C \\
\hline
$0$ & $0$  & $0$ \\
$0$ & $1$  & $0$ \\
$1$ & $0$  & $1$ \\
$1$ & $1$  & $1$ \\
\hline
\end{tabular}
\end{center}

The above table shows that the output solely depends on the 
input A. Therefore, the simplified logic is $C = A$.

\end{solutionblock}
%\fi

%%%%%%%%%%%%%%%%%%%%%%%%%%%%%%%%%%%%%%%%%%%%%%%%%%%%%%%%%%%%%
%% Task 5.2 %%
%%%%%%%%%%%%%%%%%%%%%%%%%%%%%%%%%%%%%%%%%%%%%%%%%%%%%%%%%%%%%

\task{Logic Analysis (2)}

Which function is fulfilled by the logic below ?
\begin{figure}[h!]
    \centering
    \includegraphics[width=0.65\textwidth]{task_5_2.pdf}
    % \caption{Logic Circuit}
\end{figure}

%\ifshowsolutions
\begin{solutionblock}

\begin{center}
\begin{tabular}{|>{\centering\arraybackslash}p{1cm}|
                >{\centering\arraybackslash}p{1cm}|
                >{\centering\arraybackslash}p{1cm}|
                >{\centering\arraybackslash}p{1cm}|}
\hline
$A$ & $B$ & $C$ & $C_{n+1}$\\
\hline
$0$ & $0$  & $0$ & 0\\
$0$ & $0$  & $1$ & 0\\
$0$ & $1$  & $0$ & 0\\
$0$ & $1$  & $1$ & 1\\
$1$ & $0$  & $0$ & 1\\
$1$ & $0$  & $1$ & 1\\
$1$ & $1$  & $0$ & 1\\
$1$ & $1$  & $1$ & 1\\
\hline
\end{tabular}
\end{center}
where $C_{n+1}$ is the output for the next cycle.\\
We have a feedback mechanism. Output $C$ is given as feedback.
We can set $C = 1$ with A and the to reset $C = 0$  we need $A$ and $B$.\\
Once $C = 1$ is set, this ciruit holds itself.\\
It is a Self-holding ciruit.
This results in the following simplified circuit.
\begin{figure}[h!]
    \centering
    \includegraphics[width=0.4\textwidth]{task_5_2_solution.pdf}
    % \caption{Logic Circuit}
\end{figure}

\end{solutionblock}
%\fi


%%%%%%%%%%%%%%%%%%%%%%%%%%%%%%%%%%%%%%%%%%%%%%%%%%%%%%%%%%%%%
%% Task 5.3 %%
%%%%%%%%%%%%%%%%%%%%%%%%%%%%%%%%%%%%%%%%%%%%%%%%%%%%%%%%%%%%%

\task{PLC Statement List (STL) Representation}

\subtask{What does the representation of
\\
- the conjunction (AND function) of the inputs I0.0, NOT I0.1 and I0.2  
and \\
- the assignment of the result of the logic operation RLO  
to the output Q0.0 in the form of a statement list (STL) of a PLC look like?}


\begin{solutionblock}

The corresponding STL code is:
\begin{center}
\begin{tabular}{>{\centering\arraybackslash}p{1.5cm}
                >{\centering\arraybackslash}p{1.5cm}}

\hline
$A$ & $I0.0$  \\
$AN$ & $I0.1$  \\
$A$ & $I0.2$  \\
$=$ & $Q0.0$  \\
\hline
\end{tabular}
\end{center}

\end{solutionblock}

\subtask{What does the representation of \\
- the disjunction (OR function) of the inputs I1.0, I1.1, NOT I1.2 and \\
- the assignment of the RLO to the output Q0.1 and to the  
Memory bit M 9.1 in the form of a statement list (STL) of a PLC look like?}

\begin{solutionblock}

The corresponding STL code is:

\begin{center}
\begin{tabular}{>{\centering\arraybackslash}p{1.5cm}
                >{\centering\arraybackslash}p{1.5cm}}

\hline
$O$ & $I1.0$  \\
$O$ & $I1.1$  \\
$ON$ & $I1.2$  \\
$=$ & $Q0.1$  \\
$=$ & $M9.1$  \\

\hline
\end{tabular}
\end{center}

\end{solutionblock}


%%%%%%%%%%%%%%%%%%%%%%%%%%%%%%%%%%%%%%%%%%%%%%%%%%%%%%%%%%%%%
%% Task 5.4 %%
%%%%%%%%%%%%%%%%%%%%%%%%%%%%%%%%%%%%%%%%%%%%%%%%%%%%%%%%%%%%%

\task{Contactor circuit}

The contactor circuit is given below:\\
Please find the PLC internal circuit in the Function Block Diagram (FBD)
which is equivalent to the contactor circuit !
\begin{figure}[h!]
    \centering
    \includegraphics[width=0.5\textwidth]{task_5_4.pdf}
    % \caption{Logic Circuit}
\end{figure}


\begin{solutionblock}
The following FBD represents the equivalent contactor circuit:\\
\begin{figure}[h!]
    \centering
    \includegraphics[width=0.5\textwidth]{task_5_4_solution_1.pdf}
    % \caption{Logic Circuit}
\end{figure}\\
Alternative Solution using SR Flip-Flops:
\begin{figure}[h!]
    \centering
    \includegraphics[width=0.5\textwidth]{task_5_4_solution_2.pdf}
    % \caption{Logic Circuit}
\end{figure}
\end{solutionblock}
%\fi


%%%%%%%%%%%%%%%%%%%%%%%%%%%%%%%%%%%%%%%%%%%%%%%%%%%%%%%%%%%%%
%% Task 5.5 %%
%%%%%%%%%%%%%%%%%%%%%%%%%%%%%%%%%%%%%%%%%%%%%%%%%%%%%%%%%%%%%

\task{LAD Representation of a Two-Staged Boolean Operation}

Please create the LAD representation of a two staged boolean operation
which complies with the following description:\\
The result of the OR operation with the inputs I 25.7 and I 26.0 is to be 
linked into an AND operation with the inputs I 10.6 und I 10.7. The result 
is to be assigned to the output Q 8.0.\\


%\ifshowsolutions
\begin{solutionblock}
\begin{figure}[h!]
    \centering
    \includegraphics[width=0.5\textwidth]{task_5_5_solution.pdf}
    % \caption{Logic Circuit}
\end{figure}
\end{solutionblock}


%%%%%%%%%%%%%%%%%%%%%%%%%%%%%%%%%%%%%%%%%%%%%%%%%%%%%%%%%%%%%
%% Task 5.6 %%
%%%%%%%%%%%%%%%%%%%%%%%%%%%%%%%%%%%%%%%%%%%%%%%%%%%%%%%%%%%%%

\task{Emulation of the Antivalence (XOR) Operation in a PLC}

How is the antivalence operation (XOR) to be emulated, if it is not  
contained in the basic operation pool of a certain PLC?

\subtask{STL}
\begin{solutionblock}

\begin{table}[h]
\centering
\begin{tabular}{|c|c|c|}
\hline
\textbf{A} & \textbf{B} & \textbf{A $\oplus$ B} \\ \hline
0 & 0 & 0 \\ \hline
0 & 1 & 1 \\ \hline
1 & 0 & 1 \\ \hline
1 & 1 & 0 \\ \hline
\end{tabular}
\caption{XOR Gate Truth Table}
\end{table}

\[ XOR = A\overline{B} + \overline{A}B \]
OR
\[ XOR = (A + B)(\overline{A} + \overline{B})\]
\begin{center}
\begin{tabular}{>{\centering\arraybackslash}p{1.5cm}
                >{\centering\arraybackslash}p{1.5cm}}

\hline
O$\left( \right.$ & \\
A  & I0.0 \\ 
AN & I0.1 \\
$\left. \right)$ & \\
O$\left( \right.$ & \\
AN & I0.0 \\ 
A  & I0.1 \\
$\left. \right)$ & \\
= & Q0.0\\
\hline
\end{tabular}
\end{center}
\end{solutionblock}

\subtask{LAD}
\begin{solutionblock}
\begin{figure}[h!]
    \centering
    \includegraphics[width=0.4\textwidth]{task_5_6_solution_1.pdf}
    % \caption{Logic Circuit}
\end{figure}
\end{solutionblock}

\subtask{FBD}
\begin{solutionblock}
\begin{figure}[h!]
    \centering
    \includegraphics[width=0.5\textwidth]{task_5_6_solution_2.pdf}
    % \caption{Logic Circuit}
\end{figure}
\end{solutionblock}


%%%%%%%%%%%%%%%%%%%%%%%%%%%%%%%%%%%%%%%%%%%%%%%%%%%%%%%%%%%%%
%% Task 5.7 %%
%%%%%%%%%%%%%%%%%%%%%%%%%%%%%%%%%%%%%%%%%%%%%%%%%%%%%%%%%%%%%

\task{Pulse Switch}

\textbf{Problem:} Substitution of a contactor circuit by a PLC
\begin{figure}[h!]
    \centering
    \includegraphics[width=0.5\textwidth]{task_5_7.pdf}
    % \caption{Logic Circuit}
\end{figure}
\bigskip

\textbf{List of Variables}

\begin{center}
\begin{tabular}{|c|c|c|}
\hline
Symbol & Absolute & Comment \\
\hline
S1 & I 1.0 & Momentary-contact switch \\
H0 & Q 2.2 & Lamp \\
K1 & M 0.0 & Memory bit \\
K2 & M 0.1 & Memory bit \\
K3 & M 0.2 & Memory bit \\
\hline
\end{tabular}
\end{center}

\bigskip

\textbf{Functional description:}
If the pushbutton S1 is pressed, the contactor K1 is switched on as long  
as the pushbutton is pressed. At the same time the contactor K3 is  
switched on and holds itself and the lamp H0 is on.\\
If the pushbutton S1 is pressed again, contactor K2 is switched on.  
Thereby contactor K3 drops and the lamp is off.\\
The procedure can be repeated arbitrarily.\\
The pushbutton is wired to the input I 1.0. Contactors are to be  
emulated by memory bits. The output Q 2.2 is directly connected to  
the lamp.\\
\bigskip

\textbf{Please create:}

\subtask{The hardware configuration}
\begin{solutionblock}
When switch S1 is pressed, the lamp H1 turns ON and remains ON. When the S1 switch
is again pressed, the lamp H1 turns OFF.
\begin{figure}[h!]
    \centering
    \includegraphics[width=0.5\textwidth]{task_5_7_solution_1.pdf}
    % \caption{Logic Circuit}
\end{figure}
\end{solutionblock}

\subtask{The LAD}
\begin{solutionblock}
\begin{figure}[h!]
    \centering
    \includegraphics[width=0.5\textwidth]{task_5_7_solution_2.pdf}
    % \caption{Logic Circuit}
\end{figure}

\bigskip
\end{solutionblock}


%%%%%%%%%%%%%%%%%%%%%%%%%%%%%%%%%%%%%%%%%%%%%%%%%%%%%%%%%%%%%
%% Task 5.8 %%
%%%%%%%%%%%%%%%%%%%%%%%%%%%%%%%%%%%%%%%%%%%%%%%%%%%%%%%%%%%%%

\task{Motor Control Circuit}

A function diagram of a motor control circuit is given:

\bigskip
\begin{figure}[h!]
    \centering
    \includegraphics[width=0.5\textwidth]{task_5_8.pdf}
    % \caption{Logic Circuit}
\end{figure}
\bigskip

\subtask{Please specify the hardware configuration when using a PLC!}
\begin{solutionblock}
\begin{figure}[h!]
\centering
\includegraphics[width=0.5\textwidth]{task_5_8_solution.pdf}
% \caption{Logic Circuit}
\end{figure}

\end{solutionblock}

\subtask{Please develop the statement list (STL)!}
\begin{solutionblock}
\begin{center}
\begin{tabular}{>{\centering\arraybackslash}p{1.5cm}
                >{\centering\arraybackslash}p{1.5cm}}

\hline
A$\left( \right.$ & \\
O  & A11 \\ 
O & E1 \\
$\left. \right)$ & \\
AN & A12 \\ 
A  & E3 \\
= & A11\\
\hline
\end{tabular}
\end{center}

\begin{center}
\begin{tabular}{>{\centering\arraybackslash}p{1.5cm}
                >{\centering\arraybackslash}p{1.5cm}}

\hline
A$\left( \right.$ & \\
O  & E2 \\ 
O & A12 \\
$\left. \right)$ & \\
AN & A11 \\ 
A  & E3 \\
= & A12\\
\hline
\end{tabular}
\end{center}

\end{solutionblock}

%%%%%%%%%%%%%%%%%%%%%%%%%%%%%%%%%%%%%%%%%%%%%%%%%%%%%%%%%%%%%
%% Task 5.9 %%
%%%%%%%%%%%%%%%%%%%%%%%%%%%%%%%%%%%%%%%%%%%%%%%%%%%%%%%%%%%%%

\task{Momentary Pulse in STL}

Please give a sequence of statements for a so called momentary pulse  
in STL.\\

(A momentary pulse is the logic ``1''-state of a boolean variable for  
exactly one program cycle as a reaction to the rising edge of another  
boolean variable.)

\begin{solutionblock}

\begin{figure}[h!]
\centering
\includegraphics[width=0.5\textwidth]{task_5_9_solution.pdf}
% \caption{Logic Circuit}
\end{figure}

\begin{center}
\begin{tabular}{>{\centering\arraybackslash}p{1.5cm}
                >{\centering\arraybackslash}p{1.5cm}
                >{\centering\arraybackslash}p{1.5cm}}

\hline
1 & A  & X \\ 
2 & AN & Z \\
3 & = & Y \\ 
4 & A  & X \\
5 & = & Z \\
\hline
\end{tabular}
\end{center}
Note: in step 4, we are not using AND operation. We are just reading the value of
input X and then in step 5 assigning it to Z.
\end{solutionblock}

%%%%%%%%%%%%%%%%%%%%%%%%%%%%%%%%%%%%%%%%%%%%%%%%%%%%%%%%%%%%%
%% Task 5.10 %%
%%%%%%%%%%%%%%%%%%%%%%%%%%%%%%%%%%%%%%%%%%%%%%%%%%%%%%%%%%%%%

\task{Feed Track Control in a Parts Delivery System}

The feed track in the parts delivery system should always contain  
enough parts to make sure that the picking and placing process can  
operate continuously without waiting for parts.\\
When the upper level switch sees a part for longer than 1 second,  
the motor of the hopper feeder should be switched off.  
It should be switched on if the lower level sensor sees no part for  
longer than one second.\\
A lamp indicates a failure state. This failure state is set in case of  
a congestion. It is reset by a pushbutton. The motor may not be on in  
case of a failure.

\begin{figure}[h!]
\centering
\includegraphics[width=0.8\textwidth]{task_5_10.pdf}
% \caption{Logic Circuit}
\end{figure}
\bigskip

\textbf{Wanted:} Logic in FBD

\begin{solutionblock}

\begin{figure}[h!]
\centering
\includegraphics[width=0.6\textwidth]{task_5_10_solution_1.pdf}
% \caption{Logic Circuit}
\end{figure}

\begin{figure}[h!]
\centering
\includegraphics[width=0.6\textwidth]{task_5_10_solution_2.pdf}
% \caption{Logic Circuit}
\end{figure}

\end{solutionblock}


%%%%%%%%%%%%%%%%%%%%%%%%%%%%%%%%%%%%%%%%%%%%%%%%%%%%%%%%%%%%%
%% Task 5.11 %%
%%%%%%%%%%%%%%%%%%%%%%%%%%%%%%%%%%%%%%%%%%%%%%%%%%%%%%%%%%%%%

\task{Manually operated press}

This method is usually applied to protect the hands of the operators  
who operate dangerous machines. The machine only starts if both  
activation push buttons are pressed nearly at the same time (within  
0.2 seconds). This has to happen using the left and the right hand.\\
The machine stops immediately if one of the pushbuttons is released.\\
Please develop a control program to operate the press cylinder while  
meeting the above safety requirements.\\

\begin{figure}[h!]
\centering
\includegraphics[width=0.6\textwidth]{task_5_11.pdf}
% \caption{Logic Circuit}
\end{figure}

\begin{solutionblock}
    \begin{figure}[h!]
\centering
\includegraphics[width=0.6\textwidth]{task_5_11_solution_1.pdf}
% \caption{Logic Circuit}
\end{figure}
The timer T5 keeps the signal T1 ON untill 200ms expires.

\begin{figure}[h!]
\centering
\includegraphics[width=0.6\textwidth]{task_5_11_solution_2.pdf}
% \caption{Logic Circuit}
\end{figure}

Note: The "brick on the button" trick won't work!\\
That is why edge triggered pulse is used.

\end{solutionblock}

%%%%%%%%%%%%%%%%%%%%%%%%%%%%%%%%%%%%%%%%%%%%%%%%%%%%%%%%%%%%%
%% Task 5.12 %%
%%%%%%%%%%%%%%%%%%%%%%%%%%%%%%%%%%%%%%%%%%%%%%%%%%%%%%%%%%%%%

\task{Bending Tool}

Metal sheets are bent to a certain shape using a bending tool. The  
bending sequence is as follows:\\
After the pushbutton S1 was pressed, the piston of Cylinder 1A moves  
out. Thereby the sheet is clamped and pre-bent. When the piston of  
Cylinder 1A is in its outer end position, the piston of Cylinder 2A  
moves out, finishes the bending of the sheet and is retracted after it  
has reached its outer end position. When the inner end position of  
Cylinder 2A is reached, the piston of Cylinder 1A is also retracted.\\
The technology sketch below shows the arrangement of the cylinders.\\

\begin{figure}[h!]
\centering
\includegraphics[width=0.6\textwidth]{task_5_12_1.pdf}
% \caption{Logic Circuit}
\end{figure}

The excitation of the solenoid valves was implemented up to now in a  
way as shown in the following circuit diagram:\\

\begin{figure}[h!]
\centering
\includegraphics[width=0.6\textwidth]{task_5_12_2.pdf}
% \caption{Logic Circuit}
\end{figure}

This contactor control shall be replaced by a software solution inside  
a PLC.\\

\begin{solutionblock}

Hardware configuration:\\
\begin{figure}[h!]
\centering
\includegraphics[width=0.6\textwidth]{task_5_12_solution_1.pdf}
% \caption{Logic Circuit}
\end{figure}

Symbol Table:\\
\begin{table}[h]
\centering
\begin{tabular}{|c|c|c|}
\hline
\textbf{Address} & \textbf{Symbol} & \textbf{Comment} \\ \hline
I0.1 & S1 & Push button: START \\ \hline
I0.2 & 1B1 & Inner-end pos cylinder 1A \\ \hline
I0.3 & 1B2 & Outer-end pos cylinder 1A \\ \hline
I0.4 & 2B1 & Inner-end pos cylinder 2A \\ \hline
I0.5 & 2B2 & Outer-end pos cylinder 2A \\ \hline
Q4.1 & 1M1 & Solenoid coil cylinder 1A \\ \hline
Q4.2 & 2M1 & Solenoid coil cylinder 2A \\ \hline
\end{tabular}
\end{table}

LAD:\\

\begin{figure}[h!]
\centering
\includegraphics[width=0.65\textwidth]{task_5_12_solution_2.pdf}
% \caption{Logic Circuit}
\end{figure}

FBD using Flip Flops:\\

\begin{figure}[h!]
\centering
\includegraphics[width=0.65\textwidth]{task_5_12_solution_3.pdf}
% \caption{Logic Circuit}
\end{figure}

\end{solutionblock}



%%%%%%%%%%%%%%%%%%%%%%%%%%%%%%%%%%%%%%%%%%%%%%%%%%%%%%%%%%%%%
%% Task 5.13 %%
%%%%%%%%%%%%%%%%%%%%%%%%%%%%%%%%%%%%%%%%%%%%%%%%%%%%%%%%%%%%%


\task{Person counter}

An Up-/Down-Counter is to be programmed. The persons in a room are to  
be counted. For this purpose, the entrance to the room is equipped with  
two light barriers which are installed in such a way that during the  
passage through the entrance first one and then both light barriers  
are interrupted. From that arrangement the counting signal is derived.  
The sketch shows the counting direction.\\

\begin{center}
    \includegraphics[width=0.5\textwidth]{task_5_13_1.pdf}
\end{center}

The light barriers provide a logic ``1'' signal if the beam is  
interrupted. The turnstile at the entrance is blocked if the room is  
occupied.\\
For resetting the counter, a pushbutton S3 is to be used. H6 is for  
room lighting. It is only on if there is someone in the room. K1  
serves for the turnstile control. The maximum number of persons in the  
room is 10.\\
From the timing sequence of the interruption of the beams of light,  
the counting direction can be determined. If light barrier A is  
interrupted first, then the counting direction is upward. In case of  
counting direction ``downward'', light barrier B is interrupted first.\\
\begin{center}
    \includegraphics[width=0.5\textwidth]{task_5_13_2.pdf}
\end{center}
Precondition for counting is that at least one light barrier (in this  
example light barrier A) is interrupted. If now the state of signal B  
changes from ``0'' to ``1'', the counter counts upward. If B changes  
from ``1'' to ``0'', the counter counts downward.\\
If only one light barrier is interrupted, counting is prohibited.\\
\begin{center}
    \includegraphics[width=0.5\textwidth]{task_5_13_3.pdf}
\end{center}
Here the precondition for counting is not fulfilled (signal state  
``1'' for A and an edge at light barrier B).\\

\begin{solutionblock}

Symbol Table:\\
\begin{table}[h]
\centering
\begin{tabular}{|c|c|c|}
\hline
\textbf{Address} & \textbf{Symbol} & \textbf{Comment} \\ \hline
Q4.5 & K1 & Lock turnstile \\ \hline
M1.0 & M1 & Auxiliary bit for rising edge \\ \hline
M2.0 & M3 & Auxiliary bit for falling edge \\ \hline
I0.1 & S1\_LB\_A & Light Barrier A \\ \hline
I0.2 & S2\_LB\_B & Light Barrier B \\ \hline
I0.3 & S3\_count\_R & Counter Reset \\ \hline
Q4.6 & K4 or H6 & Room Lighting \\ \hline
MW6 & MW6 & Count Value \\ \hline
\end{tabular}
\end{table}

\bigskip
Network 1: Counter 1\\
CU: Count Up when the \\
-room is not full ($\overline{K1}$)\\
-Light Barrier A is interrupted.\\
-Rising Edge of B is detected.\\

CD: Count Down when the \\
-Light Barrier A is interrupted.\\
-Falling Edge of B is detected.\\

CV: Count Value\\ 

P block: Positive Edge Detection\\
N block: Negative Edge Detection\\


\begin{center}
    \includegraphics[width=0.8\textwidth]{task_5_13_solution_1.pdf}
\end{center}

\bigskip
Network 2: Turnstile\\
$lock \rightarrow ON$ (1)\\
$unlock \rightarrow OFF$ (0)\\
\begin{center}
    \includegraphics[width=0.4\textwidth]{task_5_13_solution_2.pdf}
\end{center}

\end{solutionblock}



%%%%%%%%%%%%%%%%%%%%%%%%%%%%%%%%%%%%%%%%%%%%%%%%%%%%%%%%%%%%%
%% Task 5.14 %%
%%%%%%%%%%%%%%%%%%%%%%%%%%%%%%%%%%%%%%%%%%%%%%%%%%%%%%%%%%%%%

\task{Generation of STL from a Timing Diagram}

For the control function displayed in a timing diagram, the STL is to  
be generated.

\subtask{a}
\begin{center}
    \includegraphics[width=0.6\textwidth]{task_5_14_a.pdf}
\end{center}

\begin{solutionblock}

\begin{table}[h]
\centering
\begin{tabular}{|c|c|c|}
\hline
\textbf{I1} & \textbf{I2} & \textbf{Q1} \\ \hline
0 & 0 & 1 \\ \hline
0 & 1 & 0 \\ \hline
1 & 0 & 0 \\ \hline
1 & 1 & 0 \\ \hline
\end{tabular}
\caption{Truth Table for Logic Function (a)}
\end{table}

\begin{center}
\begin{tabular}{>{\centering\arraybackslash}p{1.5cm}
                >{\centering\arraybackslash}p{1.5cm}}

\hline
O  & I1 \\ 
O & I2 \\
=N & Q1 \\ 
\hline
\end{tabular}
\end{center}

OR

\begin{center}
\begin{tabular}{>{\centering\arraybackslash}p{1.5cm}
                >{\centering\arraybackslash}p{1.5cm}}

\hline
AN  & I1 \\ 
AN & I2 \\
= & Q1 \\ 
\hline
\end{tabular}
\end{center}

\end{solutionblock}

\subtask{b}
\begin{center}
    \includegraphics[width=0.6\textwidth]{task_5_14_b.pdf}
\end{center}
\begin{solutionblock}
    
\begin{table}[h]
\centering
\begin{tabular}{|c|c|c|}
\hline
\textbf{I1} & \textbf{I2} & \textbf{Q1} \\ \hline
0 & 0 & 0 \\ \hline
0 & 1 & 1 \\ \hline
1 & 0 & 1 \\ \hline
1 & 1 & 0 \\ \hline
\end{tabular}
\caption{Truth Table for Logic Function (b)}
\end{table}

\begin{center}
\begin{tabular}{>{\centering\arraybackslash}p{1.5cm}
                >{\centering\arraybackslash}p{1.5cm}}

\hline
O$\left( \right.$ & \\
A  & I1 \\ 
AN & I2 \\
$\left. \right)$ & \\
O$\left( \right.$ & \\
AN & I1 \\ 
A  & I2 \\
$\left. \right)$ & \\
= & Q1\\
\hline
\end{tabular}
\end{center}

\end{solutionblock}



%%%%%%%%%%%%%%%%%%%%%%%%%%%%%%%%%%%%%%%%%%%%%%%%%%%%%%%%%%%%%
%% Task 5.15 %%
%%%%%%%%%%%%%%%%%%%%%%%%%%%%%%%%%%%%%%%%%%%%%%%%%%%%%%%%%%%%%


\task{Analysis of a Statement List (STL)}

The statement list below is given in STL:

\begin{verbatim}
A     I0.0
AN    M5.0
S     M5.0
AN    I0.0
R     M5.0
\end{verbatim}

What happens here?  
Can this be done simpler?

\begin{solutionblock}

\begin{itemize}
\item If the input I0.0 is \textbf{TRUE} and the memory bit M5.0 is
\textbf{FALSE}, then M5.0 is \textbf{set}.
\item If the input I0.0 is \textbf{FALSE}, then M5.0 is
\textbf{reset}.
\end{itemize}

Thus, the memory bit M5.0 follows the state of the input I0.0.  
M5.0 is set when I0.0 becomes TRUE and reset when I0.0 becomes FALSE.\\

This behavior corresponds to a simple assignment of the input I0.0 to
the memory bit M5.0.\\

The same function can be implemented in a simpler and clearer way using
a direct assignment:\\

\begin{center}
\begin{tabular}{>{\centering\arraybackslash}p{1.5cm}
                >{\centering\arraybackslash}p{1.5cm}}

\hline
A  & I0.0 \\ 
= & M5.0\\
\hline
\end{tabular}
\end{center}
\end{solutionblock}
